\begin{thebibliography}{99}
\footnotesize\itemsep=0pt

\bibitem{BSV}
Bozec T., Schiffmann O., Vasserot E., On the number of points of nilpotent
 quiver varieties over finite fields,
 \href{https://doi.org/10.24033/asens.2452}{\textit{Ann. Sci. \'Ec. Norm.
 Sup\'er.}} \textbf{53} (2020), 1501--1544,
 \href{http://arxiv.org/abs/1701.01797}{arXiv:1701.01797}.

\bibitem{B}
Bridgeland T., An introduction to motivic {H}all algebras,
 \href{https://doi.org/10.1016/j.aim.2011.09.003}{\textit{Adv. Math.}}
 \textbf{229} (2012), 102--138,
 \href{http://arxiv.org/abs/1002.4372}{arXiv:1002.4372}.

\bibitem{E}
Efimov A.I., Cohomological {H}all algebra of a symmetric quiver,
 \href{https://doi.org/10.1112/S0010437X12000152}{\textit{Compos. Math.}}
 \textbf{148} (2012), 1133--1146,
 \href{http://arxiv.org/abs/1103.2736}{arXiv:1103.2736}.

\bibitem{FH}
Fine N.J., Herstein I.N., The probability that a matrix be nilpotent,
 \href{https://doi.org/10.1215/ijm/1255454112}{\textit{Illinois~J. Math.}}
 \textbf{2} (1958), 499--504.

\bibitem{Hesse}
Hesselink W.H., Desingularizations of varieties of nullforms,
 \href{https://doi.org/10.1007/BF01390087}{\textit{Invent. Math.}} \textbf{55}
 (1979), 141--163.

\bibitem{H}
Hilbert D., Ueber die vollen {I}nvariantensysteme,
 \href{https://doi.org/10.1007/BF01444162}{\textit{Math. Ann.}} \textbf{42}
 (1893), 313--373.

\bibitem{KW}
Kraft H., Wallach N.R., On the nullcone of representations of reductive groups,
 \href{https://doi.org/10.2140/pjm.2006.224.119}{\textit{Pacific~J. Math.}}
 \textbf{224} (2006), 119--139.

\bibitem{LB}
Le~Bruyn L., Optimal filtrations on representations of finite-dimensional
 algebras, \href{https://doi.org/10.1090/S0002-9947-00-02590-3}{\textit{Trans.
 Amer. Math. Soc.}} \textbf{353} (2001), 411--426.

\bibitem{LBB}
Le~Bruyn L., Noncommutative geometry and {C}ayley-smooth orders, \textit{Pure
 Appl. Math.}, Vol. 290, Chapman \& Hall/CRC, Boca Raton, FL, 2008.

\bibitem{LP}
Le~Bruyn L., Procesi C., Semisimple representations of quivers,
 \href{https://doi.org/10.2307/2001477}{\textit{Trans. Amer. Math. Soc.}}
 \textbf{317} (1990), 585--598.

\bibitem{MozIntro}
Mozgovoy S., Introduction to {D}onaldson--{T}homas invariants, in Advances in
 {R}epresentation {T}heory of {A}lgebras, \textit{EMS Ser. Congr. Rep.}, European
 Mathematical Society, Z\"urich, 2013, 195--210.

\bibitem{RCounting}
Reineke M., Counting rational points of quiver moduli,
 \href{https://doi.org/10.1155/IMRN/2006/70456}{\textit{Int. Math. Res. Not.}}
 \textbf{2006} (2006), 70456, 19~pages,
 \href{http://arxiv.org/abs/math.AG/0505389}{arXiv:math.AG/0505389}.

\bibitem{RSmall}
Reineke M., Quiver moduli and small desingularizations of some {GIT} quotients,
 in Representation {T}heory~-- {C}urrent {T}rends and {P}erspectives, \textit{EMS Ser.
 Congr. Rep.}, \href{https://doi.org/10.4171/171-1/21}{European Mathematical
 Society}, Z\"urich, 2017, 613--635,
 \href{http://arxiv.org/abs/1511.08316}{arXiv:1511.08316}.

\bibitem{RDT}
Reineke M., Donaldson--{T}homas invariants of symmetric quivers,
 \href{http://arxiv.org/abs/2410.03219}{arXiv:2410.03219}.

\end{thebibliography}
